\chapter{Process Flow: Reconcile and Temporal Correctness}
Reconciliation is what separates a real memory from a pile of embeddings. It is the moment a new
fact meets an old one.

\section{The Decision Table}
\textbf{Trigger:} An embedded candidate arrives from the Scribe.

\textbf{System Logic:} The reconciler searches live memories in the same owner, kind, and
quarantine partition (quarantined candidates match quarantined rows only, so poison can never
supersede a user fact) for \texttt{(subject, predicate)} equality \emph{or} cosine above
\texttt{TOTTIE\_RECONCILE\_COSINE} (0.86) with an anchor (shared tag or SP-partial; mutual
\texttt{NULL}s are not an anchor). Event-kind candidates skip matching (episodes accumulate).
Match rows are \texttt{FOR UPDATE} locked inside the turn's transaction; the per-partition live-SP
unique index backstops races and the Scribe retries the turn on violation (verified: four
concurrent twins converge to one live row with a 3-link chain). Every \texttt{NEW} also writes
\texttt{tag} links (capped at 3, same partition), and every \texttt{SUPERSEDE} writes a
\texttt{supersedes} link, so the graph substrate grows with the store. Then it decides:

\begin{longtable}{@{}p{0.16\textwidth}p{0.34\textwidth}p{0.42\textwidth}@{}}
\toprule
Operation & When & What happens \\
\midrule
\texttt{NEW} & No prior memory matches & Insert with \texttt{valid\_from = now()} \\
\texttt{SKIP} & A matching memory already holds the same value & Bump confidence / salience; no new row \\
\texttt{UPDATE} & A matching memory is enriched but not contradicted & Append detail; refresh summary \\
\texttt{SUPERSEDE} & A matching memory holds a \emph{different} value & Set old \texttt{valid\_to = now()}; insert new with \texttt{supersedes = old.id} \\
\bottomrule
\end{longtable}

\textbf{Data Flow \& Mutators:} Any insert or update is written to \texttt{memories}; the decision is
written to \texttt{memory\_event.scribe\_action} and streamed to the Panel, where a superseded memory
is rendered with a strikethrough.

\section{Worked Example: ``I'm vegetarian now''}
\textbf{Trigger:} The user says \emph{``actually, I'm vegetarian now''}.

\textbf{Action / UI State:} The Panel's ``MEMORIES FORMED'' block shows a red-ish
\texttt{SUPERSEDE} badge, the old value struck through (``diet: meat-eater''), and the new value
(``diet: vegetarian'').

\textbf{System Logic:}
\begin{enumerate}
  \item The Observer flags a durable preference change (salience high; \texttt).
  \item The Scribe embeds the candidate and finds the live fact \texttt{(subject=user,
  predicate=diet, object=meat-eater)}.
  \item The values differ, so the operation is \texttt{SUPERSEDE}.
  \item The old row's \texttt{valid\_to} is set to \texttt{now()}; the new row is inserted with
  \texttt{valid\_from = now()} and \texttt{supersedes = old.id}.
\end{enumerate}

\textbf{Data Flow \& Mutators:} Two row mutations and one event; retrieval now returns only
``vegetarian'' because it filters \texttt{valid\_to IS NULL}. The agent never has to reconcile
in-prompt, and it can never quote a stale fact.

\section{Why This Is Hard and Why It Wins}
Plain vector recall has \emph{no update semantics}: both the old and the new fact retrieve, and the
model is left to guess which is current. Tottie resolves the conflict at write time, once, and
records the temporal chain. This is precisely the capability that LongMemEval's \emph{knowledge
update} category and MemoryAgentBench's \emph{selective forgetting / conflict resolution} category
probe --- and it is a headline slide in the demo.
